\originalchapter{8.334作业4：转移矩阵与实空间重整化}
\psource{Mehran Kardar，8.334 Spring 2014，PS4，截止2014年4月9日。原卷前四页介绍转移矩阵，理论第\ref{ch:duality}章及本章承接，全部七题包括选做题收入。}
\begin{exercise}[原题1(a)--(d)]\label{p:334-4-1}
一维spin1取$s=-1,0,1$，$-\beta H=K\sum s_is_{i+1}$。(a)转移矩阵；(b)利用对称求最大本征值与$\log Z/N$；(c)相关长和$K\to\infty$；(d)消去隔站后闭合的最简参数空间。
\end{exercise}
\begin{solution}
按$(-1,0,1)$，$T=\left(\begin{smallmatrix}\e^K&1&\e^{-K}\\1&1&1\\\e^{-K}&1&\e^K\end{smallmatrix}\right)$，奇向量$(1,0,-1)/\sqrt2$的$\lambda_o=2\sinh K$。偶空间取$(1,0,1)/\sqrt2,(0,1,0)$，块为$\left(\begin{smallmatrix}2\cosh K&\sqrt2\\\sqrt2&1\end{smallmatrix}\right)$，最大
\begin{equation}
\lambda_+=\frac{2\cosh K+1+\sqrt{(2\cosh K-1)^2+8}}2,\qquad \log Z/N\to\log\lambda_+.
\end{equation}
自旋关联连接偶基态与奇模，$\xi^{-1}=\log(\lambda_+/|2\sinh K|)$；$K\to+\infty$时$\lambda_+-\lambda_o\simeq4\e^{-K}$，$\xi\simeq\e^{2K}/4$。有限温正矩阵最大根单一，无非零温相变。

闭合核可写$g+Jss'+Qs^2s'^2-D(s^2+s'^2)/2$。四个独立对称权重$T_{00},T_{01},T_{11},T_{1,-1}$都可拟合这四参数，平方$T^2$保持全局$s\mapsto-s$对称，却生成晶场$D$与双二次$Q$，仅$J$不能闭合。
\end{solution}
\begin{exercise}[原题2(a)--(d)]\label{p:334-4-2}
$q$态clock自旋$s\in\mathbb Z_q$，近邻$-\beta H=\sum J(s_i-s_j)$，$J$为任意周期实函数。(a)环链$Z=\Tr T^N$；(b)显写循环矩阵并离散Fourier对角化；(c)热力学自由能；(d)以投影$\Pi_\alpha$写$\avg{\delta_{s_i,s_{i+l}}}$，原题声称连接部分$\sim[\lambda(1)/\lambda(0)]^l$，说明条件。
\end{exercise}
\begin{solution}
$w_n=\e^{J(n)}$、$T_{ab}=w_{a-b}$，环乘积求和即迹。Fourier本征向量$q^{-1/2}\e^{-2\pi\ii ka/q}$，$\lambda_k=\sum_{n=0}^{q-1}w_n\e^{2\pi\ii kn/q}$，$
Z=\sum_k\lambda_k^N$。严格正权重给Perron根$\lambda_0=\sum w_n>|\lambda_{k\ne0}|$，$\beta f=-\log\lambda_0$。

$\Pi_\alpha$是对角仅第$\alpha$为1，$\avg{\delta_{s_i,s_{i+l}}}=Z^{-1}\sum_\alpha\Tr(\Pi_\alpha T^l\Pi_\alpha T^{N-l})$。无穷链由均匀稳态得
\begin{equation}
\avg{\delta_{s_i,s_{i+l}}}_c=\frac1q\sum_{k=1}^{q-1}(\lambda_k/\lambda_0)^l.
\end{equation}
一般$J(n)$下主导为非零$k$中的最大模，未必$k=1$；例如$q=4$、$w_0=w_2>w_1=w_3$，$\lambda_2$主导而$\lambda_1=0$。原题末式对通常铁磁cosine clock中$k=\pm1$最慢模成立，但不能要求证明任意周期$J$下的错误断言。
\end{solution}
\begin{exercise}[原题3(a)--(c)]\label{p:334-4-3}
一维XY单位自旋$-\beta H=K\sum\cos(\theta_i-\theta_{i+1})$。(a)转移核与$\e^{\ii m\theta}$本征函数；(b)自由能及$T\propto K^{-1}\to0$热容；(c)相关長与$K\to\infty$。
\end{exercise}
\begin{solution}
用$\dd\theta$测度，$T(\theta,\theta')=\e^{K\cos(\theta-\theta')}$，卷积本征值$\lambda_m=2\pi I_m(K)$，$f=-\kb T\log[2\pi I_0(K)]$。$K=\widehat J/(\kb T)$，能每站$-\widehat J I_1/I_0$，$I_0(K)\simeq\e^K/\sqrt{2\pi K}$给$U/N=-\widehat J+\kb T/2+\cdots$，所以经典自旋$C/N\to\kb/2$，这是独立相对角谐振模；若量子自由度会有低温量子修正。

$s_i\cdot s_{i+l}$关联取$m=\pm1$，$\xi^{-1}=\log(I_0/I_1)$，$I_1/I_0=1-1/(2K)+O(K^{-2})$，$\xi\simeq2K$。与离散Ising的激活式$\e^{2K}$不同，连续角涨落用渐进的小角变化破坏一维有序。
\end{solution}
\begin{exercise}[原题4(a)--(f)，选做]\label{p:334-4-4}
一维经典气有足够硬核使只有最近粒子相互作用，排序$0\le x_1\le\cdots\le x_N\le L$。(a)写$H=\sum p_i^2/(2m)+\sum_{i=2}^NV(x_i-x_{i-1})$；(b)$Z$及$\delta_1=x_1,\delta_i=x_i-x_{i-1}$区域；(c)Gibbs--Laplace配分$\mathcal Z(P)=\int_0^\infty\e^{-\beta PL}Z(L)\dd L$与鞍点压力；(d)$\delta_{N+1}=L-\sum_1^N\delta_i$分离积分；(e)平均长、密度与无凝聚；(f)硬棒最小间距$a$的状态方程及virial系数。
\end{exercise}
\begin{solution}
排序去掉$N!$，动量积分后$Z=\lambda^{-N}\int_{\delta_i\ge0,\sum_1^N\delta_i\le L}\prod_{i=1}^N\dd\delta_i\exp[-\beta\sum_{i=2}^NV(\delta_i)]$。两端外侧没有粒子间势。Laplace鞍点$\partial_L\log Z=\beta P$，就是canonical压力。

最后空隙变量使$L=\sum_1^{N+1}\delta_i$，Jacobian1。两个边界空隙积分各$(\beta P)^{-1}$，$I(P,T)=\int_0^\infty\e^{-\beta[V(\delta)+P\delta]}\dd\delta$，故$\mathcal Z=\lambda^{-N}(\beta P)^{-2}I^{N-1}$。$\avg L=-\beta^{-1}\partial_P\log\mathcal Z=2/(\beta P)+(N-1)\avg\delta$，$\avg\delta=\int\delta\e^{-\beta(V+P\delta)}/I$，热力学$n=1/\avg\delta$。对稳定短程势、正压有限温，正积分解析，无均场所预言的有限温液气奇点；不是对任意不稳定势都成立。

硬棒$I=\e^{-\beta Pa}/(\beta P)$，$\avg\delta=a+\kb T/P$，故$P=n\kb T/(1-na)$，$B_l=a^{l-1}$。有限$N$的边界空隙给$\avg L=(N-1)a+(N+1)\kb T/P$，不能误当无穷体系方程。
\end{solution}
\begin{exercise}[原题5(a)--(i)]\label{p:334-4-5}
一维$q$态Potts环链$-\beta H=J\sum\delta_{s_i,s_{i+1}}$。(a)$Z,\xi$；(b)反铁磁零温是否临界；(c)隔站消去的$J,g$流；(d)固定点稳定；(e)$\log Z$的常数级数；(f)$t=\e^{-1/\xi}$化简；(g)级数回到转移矩阵答案；(h)加入$h\sum\delta_{s_i,1}$生成$K\sum\delta_{s_i,1}\delta_{s_{i+1},1}$的三参数流；(i)铁磁零温的磁本征值与相关长。
\end{exercise}
\begin{solution}
令$x=\e^J$，$T_{ab}=1+(x-1)\delta_{ab}$，$\lambda_0=x+q-1$、$\lambda_1=x-1$（$q-1$重），$Z=\lambda_0^N+(q-1)\lambda_1^N$、$\xi^{-1}=\log(\lambda_0/|\lambda_1|)$。$J\to-\infty$时$|\lambda_1/\lambda_0|=1/(q-1)$，$q>2$相关长有限，$q=2$为严格交替零温临界，不能统一回答所有$q$。

$T^2$非对角$D=2x+q-2$、对角$C=x^2+q-1$，所以$J'=\log(C/D)$、$g'=2g+\log D$。$J=0$稳定，铁磁$J=\infty$零温点不稳定；对$q=2$反铁磁极限消去隔站后映到铁磁点，交错序需另外跟踪。

铁磁取$t=(x-1)/(x+q-1)\ge0$，$t'=t^2$；反铁磁用有符号比$r$才有$r'=r^2$。从$g_0=0$，每原站$\log Z/N=\sum_{j\ge0}2^{-j-1}\log D(J_j)$。$D=q(1+t)/(1-t)$，因此级数$\log q+\sum2^{-j-1}\log[(1+t^{2^j})/(1-t^{2^j})]$，由$1-t^{2^{j+1}}=(1-t^{2^j})(1+t^{2^j})$逐项抵消为$\log[q/(1-t)]$，正是$\log\lambda_0$。

三参数最清楚从对称核$T_{ab}=\exp[g+J\delta_{ab}+K\delta_{a1}\delta_{b1}+h(\delta_{a1}+\delta_{b1})/2]$平方。按$11$、$1a$、$aa$、$ab$（$a,b\ne1,a\ne b$）四权重去掉$\e^{2g}$后分别
\begin{align}
A&=\e^{2(J+K+h)}+(q-1)\e^h,& B&=\e^{h/2}[\e^{J+K+h}+\e^J+q-2],\\
C&=\e^h+\e^{2J}+q-2,&D&=\e^h+2\e^J+q-3.
\end{align}
故$g'=2g+\log D$、$J'=\log(C/D)$、$h'=2\log(B/D)$、$K'=\log(A/D)-J'-h'$。$q=2$的参数存在冗余，可按Ising场直接闭合；$q\ge3$上式四类齐备。

$J\to\infty$有$J'=J-\log2+\cdots$，$\e^{-J}$翻倍，热本征值2、$y_t=1$；$\xi\simeq\e^J/q$。磁小扰动$\binom{h'}{K'}=\left(\begin{smallmatrix}2&1\\0&1\end{smallmatrix}\right)\binom hK+\cdots$，本征值2、1，对应$y_h=1$与额外边缘方向0。有限温一维没有自发破缺，零温磁场敏感不能推出非零$T_c$。
\end{solution}
\begin{exercise}[原题6(a)--(i)]\label{p:334-4-6}
蜂窝Ising近邻$J$。(a)四站星形簇、长度缩放$b=2$；(b)多数规则如何处理平票；(c)中心为平票决胜的簇自旋$+$微观表；(d)相邻簇内部与跨边能；(e)令$x=\e^J$，证明$R_{++}=x^8+2x^6+7x^4+14x^2+17+14x^{-2}+7x^{-4}+2x^{-6}$，$R_{+-}=9x^4+16x^2+13+16x^{-2}+9x^{-4}+x^{-8}$；(f)耦合流；(g)近似临界$J_c,\nu$与精确值比较；(h)零温铁磁的$y_h,y_t$；(i)反铁磁是否适用及破坏的对称。
\end{exercise}
\begin{solution}
星形取中心$s_0$与三个叶$s_1,s_2,s_3$，总和非0用多数符号，等于0用$s_0$。蜂窝适当四站覆盖后两相邻星形有两条跨簇边。$+$微观态8个：中心$+$、叶中负数0、1、2（分别1、3、3种）；中心$-$仅叶全$+$（1种）。$-$态为全体反号，表已穷尽16态。

设两个簇$a,b$，可标跨边$a_1b_1,a_2b_2$，总键和$E=\sum_{j=1}^3(a_0a_j+b_0b_j)+a_1b_1+a_2b_2$，权重$\e^{JE}$。固定簇$+,+$或$+,-$各8乘8种，按$E$计数得到题中两多项式；每式在$x=1$均64，最高键和8验证总8条微观键。随卷数值代码完整穷举验证系数。

局部两簇投影$R(S_A,S_B)=A_0\exp[J'S_AS_B]$，故$J'=\tfrac12\log(R_{++}/R_{+-})$。这一局部权重近似忽略相邻第三簇、多簇有效相互作用，不是整个格点分区函数的精确闭合变换。非零固定点数值$J_c\simeq1.04553$，$\lambda_t=J'(J_c)'\simeq1.81762$，$\nu=\log2/\log\lambda_t\simeq1.1600$；精确蜂窝$J_c=\tfrac12\log(2+\sqrt3)=0.65848$、$\nu=1$。较大非普适误差必须如实报告，不调系数“校准”到精确答案。

铁磁零温$J'\simeq2J-\log3$，以$T\propto1/J$坐标$T'\simeq T/2$，$y_t=-1$；一个饱和四站簇磁功$4hS$，$y_h=2$。若用激活缺陷$\e^{-2J}$作坐标，线性导数为0，与$1/J$坐标不同，不能伪称有限激活本征值。反铁磁Néel态中心与三个叶反号，簇多数投影与中心决胜不与交错翻转规范变换相容，不能保留原蜂窝两亚格子对称；需改交错块变量后另建方案。
\end{solution}
\begin{exercise}[原题7(a),(b)，选做]\label{p:334-4-7}
两个Ising自旋$H=-\widehat K\sigma_1\sigma_2$，$K=\beta\widehat K$。(a)均匀随机选一个自旋做Metropolis，显写四态转移矩阵并验证平衡权重为左本征向量；(b)单翻、双翻以$p,1-p$选取，重复。
\end{exercise}
\begin{solution}
顺序$(++,+-,-+,--)$，令$a=\min(1,\e^{-2K})$、$b=\min(1,\e^{2K})$。行代表初态，
\begin{equation}
P_1=\begin{pmatrix}1-a&a/2&a/2&0\\b/2&1-b&0&b/2\\b/2&0&1-b&b/2\\0&a/2&a/2&1-a\end{pmatrix},\quad
P_2=\begin{pmatrix}0&0&0&1\\0&0&1&0\\0&1&0&0\\1&0&0&0\end{pmatrix}.
\end{equation}
$\pi=(\e^K,\e^{-K},\e^{-K},\e^K)/[4\cosh K]$，$\e^Ka=\e^{-K}b$使每边详细平衡，直接得$\pi P_1=\pi$。双翻不改键能，必接受，且配对权重相等，$\pi P_2=\pi$，故$P=pP_1+(1-p)P_2$同样平稳。$p=0$分成两不连通轨道，虽然平稳式仍真，却不能从任意初态采到完整平衡；平稳、不可约、有效混合是三个不同问题。
\end{solution}
